\chapter{Related Work}
\label{cha:related}

This chapter surveys prior work across five research areas directly relevant to the contributions of this thesis: tiny and edge-deployable vision-language models, model quantization for efficient inference, episodic memory in language and multimodal systems, hallucination detection and control in VLMs, and mission-critical visual recognition benchmarks. Each section contextualizes prior work relative to the thesis contributions and concludes with a statement of how this thesis advances beyond the existing literature.

\section[Tiny and Edge-Deployable VLMs]{Tiny and Edge-Deployable Vision-Language Models}
\label{sec:related_tiny_vlm}

The rapid scaling of VLMs has prompted a parallel line of work on compact architectures that preserve multimodal capability under reduced parameter budgets. SmolVLM~\cite{smolvlm2024} demonstrates that sub-billion-parameter models paired with strong vision encoders can achieve competitive results on standard benchmarks while being deployable on commodity hardware. The authors show that careful data curation and efficient vision tokenization significantly impact quality at small scales. However, even SmolVLM-256M requires a discrete GPU for practical inference and does not target CPU-only deployment.

MobileVLM~\cite{chu2023mobilevlm} proposes lightweight vision-language connectors and mobile-optimized language backbones for responsive on-device inference. The architecture demonstrates that the vision-language bridge is a critical bottleneck in compact VLMs, and that lightweight depthwise-separable designs can reduce parameter count. MobileVLM V2~\cite{chu2024mobilevlmv2} extends this with token compression and improved throughput. While these models represent genuine progress, they target GPU-based mobile devices rather than the CPU-only, RAM-limited deployments addressed in this thesis.

LLaVA-Phi~\cite{zhu2024llavaphi} combines a compact Phi-2 language model with a CLIP vision encoder, showing that 2.7B-parameter instruction-tuned VLMs can approach much larger models at substantially lower inference cost. This is relevant to robotics and edge deployment, but 2.7B parameters still exceeds the practical limits of the target hardware by more than an order of magnitude. TinyGPT-V~\cite{yuan2024tinygptv} further reduces parameters by employing efficient patch embedding strategies and distillation from larger teacher models. MiniCPM-V~\cite{hu2024minicpmv} pursues efficiency through selective parameter sharing and dynamic token pruning, achieving strong results in the 2--3B parameter range. H2OVL-Mississippi targets inference on low-resource hardware by combining aggressive pruning with knowledge distillation.

The consistent pattern across all prior compact VLM work is that none simultaneously targets CPU-only inference, under 800~MB RAM, sub-10ms per-sample latency, and zero-hallucination reliability for mission-critical classification. Most assume GPU acceleration; all treat hallucination control and epistemic memory as orthogonal concerns rather than core design requirements. This thesis addresses the full intersection of these requirements with a unified architecture.

\section{Model Quantization for Efficient Inference}
\label{sec:related_quant}

Quantization reduces model memory and compute requirements by representing weights and/or activations at lower bit-widths than standard 32-bit floating point. BitNet~\cite{wang2023bitnet} demonstrates that transformer language models can be trained from scratch with 1-bit weights without catastrophic accuracy degradation, enabling highly memory-efficient inference on specialized hardware. The key insight is that signed integer multiply-accumulate operations can replace expensive floating-point operations throughout the forward pass.

BitNet~b1.58~\cite{ma2024bitnet158} extends this to ternary weights $\{-1, 0, +1\}$, which require only 1.58 bits per parameter (the binary entropy of a ternary distribution). At this precision, nearly all multiply operations can be replaced by additions and subtractions, dramatically reducing compute on hardware without dedicated low-precision units. BitNet~b1.58 achieves near-lossless perplexity compared to full-precision baselines when trained at sufficient scale, validating ternary quantization as a practical compression strategy.

GPTQ~\cite{frantar2022gptq} provides one-shot post-training quantization via approximate second-order information, enabling efficient 3--4 bit quantization of already-trained large models without backpropagation. This is complementary to BitNet-style native quantization and is widely used to compress pre-trained models for deployment. AWQ~\cite{lin2023awq} proposes activation-aware weight quantization that identifies and protects salient channels most sensitive to quantization error, achieving strong 4-bit compression across language and multimodal models.

BitMar, the first contribution of this thesis, applies 1.58-bit ternary quantization to a multimodal pipeline: both the text encoder and all downstream fusion, memory, and decoder projections operate under ternary weight constraints. This is a direct extension of BitNet~b1.58's quantization scheme to the cross-modal setting, and to the best of our knowledge is the first architecture to demonstrate that ternary quantization is compatible with cross-modal attention fusion and episodic memory conditioning simultaneously.

\section[Episodic Memory in Neural Models]{Episodic Memory in Language and Multimodal Models}
\label{sec:related_memory}

The idea of augmenting neural networks with external, addressable memory has a long history. Differentiable Neural Computers (DNC)~\cite{graves2016dnc} introduced the first complete architecture for learning explicit read and write operations over an external memory matrix, demonstrating that neural networks can solve complex algorithmic tasks requiring long-range structured recall. While DNCs operate only on symbolic sequences, they established the key design principles, including content-based addressing, soft writes, and read-weight computation, that inform modern memory-augmented VLMs.

Memorizing Transformers~\cite{wu2022memorising} show that augmenting standard transformer language models with external memory lookups over past key-value pairs improves long-context modeling without changing core model weights. This work demonstrates that memory retrieval can be separated from parametric knowledge, allowing a frozen model to leverage growing context stores. Retrieval-Augmented Generation (RAG)~\cite{lewis2020rag} takes a related approach, connecting generative models to large external document stores to improve factuality and reduce hallucination on knowledge-intensive tasks. However, RAG targets server-scale infrastructure: dense passage retrieval from millions of documents requires GPU-accelerated vector search engines that are impractical at the edge.

Larimar~\cite{das2024larimar} presents a memory-controlled LLM framework in which episodes are written to and retrieved from a structured store using a differentiable interface, enabling the model to update its knowledge without full retraining. This is conceptually close to the EmberVLM episodic memory design, though Larimar targets large language models on GPU servers rather than tiny VLMs on CPU-only edge hardware.

EGO and related selective episodic memory strategies explore how memory writes can be conditioned on novelty signals to prevent the memory matrix from being overwritten with redundant common experiences. This informs EmberVLM's novelty-triggered write mechanism, where a minimum-distance novelty score gates whether an incoming embedding warrants a new memory write.

A critical distinction between the memory systems in this thesis and server-scale RAG or database-backed retrieval: both BitMar and EmberVLM use \emph{fixed-size, on-device episodic matrices} that fit entirely within the primary compute device's RAM (1.2~MB). No external database query, network request, or dense retrieval index is involved. Memory reads and writes are pure matrix operations that complete in microseconds. This makes the memory system viable for real-time edge inference in a way that server-scale retrieval is not.

\section{Hallucination Detection and Control in VLMs}
\label{sec:related_hallucination}

Hallucination in VLMs, the generation of tokens unsupported by the input image, has been extensively documented and benchmarked. Rohrbach~et~al.~\cite{rohrbach2018hallucination} showed that image captioning models systematically generate statistically likely but visually absent objects, even in high-performing systems. They demonstrate that hallucination correlates with co-occurrence statistics in training data rather than actual visual evidence. This foundational finding motivates all subsequent hallucination mitigation research.

The POPE benchmark~\cite{li2023pope} formalizes object hallucination evaluation as a polling-based yes/no question protocol: given an image, models are asked whether specific objects are present. POPE includes adversarial splits where absent objects are specifically chosen to have high co-occurrence with genuinely present objects, making them maximally likely to be hallucinated. POPE has become the standard evaluation for object-level hallucination across the VLM literature.

Visual Contrastive Decoding (VCD)~\cite{leng2024vcd} addresses hallucination without retraining by contrasting output distributions from original and spatially distorted visual inputs. Tokens whose generation probability is not significantly affected by the visual content are penalized, encouraging the model to rely on actual image features rather than language priors. VCD is effective across multiple large VLM families, but it requires a complete second forward pass with perturbed input, doubling the inference cost. For the CPU-only, sub-10ms deployment targets of this thesis, this overhead is unacceptable.

Instruction-following calibration approaches train models to express uncertainty when visual evidence is insufficient, requiring either additional labeled data or specialized fine-tuning pipelines. Grounding supervision methods add explicit visual localization objectives during training to reduce the gap between language-prior generation and visually evidenced generation. Both classes of methods require retraining or supplementary annotation, making them infeasible for deployment-time use on frozen edge models.

The Visual Absence (VA) Refiner introduced in this thesis occupies a previously unexplored position in the mitigation landscape: it is training-free, requires no extra forward passes, needs no labeled grounding data, and operates by monitoring the model's own internal neuron-level activation patterns during a single standard forward pass. By hooking into mid-decoder FFN layers (10--14 of SmolLM-135M), the VA Refiner detects the characteristic activation signature of hallucinated tokens, specifically symmetric patterns that do not change when the visual input is removed, and applies type-aware logit penalties to suppress them. This mechanistic grounding distinguishes the VA Refiner from all prior approaches and makes it uniquely suitable for frozen edge deployment.

\section{Mission-Critical Visual Recognition Benchmarks}
\label{sec:related_benchmarks}

Standard VLM benchmarks such as VQA v2~\cite{VQA}, GQA~\cite{hudson2019gqa}, and SugarCrepe~\cite{hsieh2023sugarcrepe} are designed to measure general visual question answering and compositional reasoning capabilities across diverse image domains. While valuable for comparative evaluation, these benchmarks do not capture the reliability requirements of safety-critical deployment. A model that achieves high GQA accuracy may still produce catastrophically wrong outputs in emergency or disaster scenarios.

The Top-N Robot Selection task, used as the primary evaluation in this thesis, is a multi-label morphology classification problem: given a scene image and mission description, predict which robot types (Drone, Underwater Robot, Humanoid, Wheeled, Legged) should be deployed. This task is mission-critical because wrong predictions directly lead to incorrect asset allocation in emergency response scenarios.

VERI-Emergency~\cite{choi2026better}, introduced at WACV 2026, is a visual emergency recognition benchmark covering fire, flooding, chemical hazards, and related emergency scenarios. It introduces an asymmetric evaluation metric structure distinguishing OverRx (false positive rate: wrongly flagging non-emergencies) from UnderRx (false negative rate: missing genuine emergency signals). For safety-critical deployment, UnderRx is the more consequential metric: a missed emergency is more dangerous than a false alarm.

GEO-Bench-VLM~\cite{danish2025geobench}, introduced at ICCV 2025, is a geospatial disaster understanding benchmark covering satellite and aerial imagery of flood zones, wildfire damage, earthquake devastation, and infrastructure collapse. It evaluates models on fine-grained scene classification tasks requiring both visual recognition and geographical contextual reasoning.

Together, these three benchmarks provide a more realistic and demanding evaluation surface for edge-deployable VLMs than standard generative benchmarks. EmberVLM achieves state-of-the-art performance across all three among sub-billion-parameter models, while simultaneously meeting strict edge-deployment latency and memory constraints.

\section{Low-Carbon and Sustainable AI}
\label{sec:related_carbon}

The environmental cost of training large AI models has attracted increasing attention. Strubell~et~al.~\cite{strubell2019energy} estimated that training a large NLP model from scratch produces carbon emissions comparable to five automobile lifetimes. More recent estimates place the training carbon footprint of GPT-3 at approximately 552 tonnes CO$_2$eq~\cite{patterson2021carbon}. As VLMs grow to hundreds of billions of parameters, their environmental impact grows correspondingly.

Compact, parameter-efficient model design is therefore not only a systems engineering concern but an environmental imperative. By constraining total trainable parameters to under 40M and completing all training stages on a dual-A100 setup in approximately 90 hours, EmberVLM achieves a total measured training footprint of 25.3~kg CO$_2$eq, verified using the CodeCarbon~\cite{codecarbon} framework. This is more than 21,000 times smaller than GPT-3's estimated footprint, demonstrating that mission-critical capable multimodal AI can be trained sustainably. This thesis treats low-carbon training not as a secondary benefit but as a primary design constraint.
